\documentclass[12pt]{article}
\usepackage[T1]{fontenc}
\usepackage[utf8]{inputenc}
\usepackage{lmodern}
\usepackage{textcomp}
\usepackage{enumitem}
\usepackage{geometry}
\geometry{margin=1in}
\begin{document}
\begin{center}
Answer Key - Mathematics - K-12 Level Assessment 25 \
\small (Answers are ordered exactly as in the question paper.)
\end{center}

\section*{Multiple Choice}
\begin{enumerate}[label=\arabic*.]
\item \textbf{Answer:} D
\\ \textit{Options:} A) \ensuremath{\frac}\{3\ensuremath{\sqrt{3}}\}\{3\}; B) \ensuremath{\frac{1}{3}}; C) \ensuremath{\sqrt{3}}; D) \ensuremath{\frac}\{\ensuremath{\sqrt{3}}\}\{3\}
\\ \textit{Note:} The expression simplifies to \(\frac{\sqrt{3}}{3}\) because \(\frac{1}{729}\) equals \(729^{-1}\), which simplifies through successive roots to yield \(\sqrt{3}/3\) after rationalizing the denominator.
\item \textbf{Answer:} B
\\ \textit{Options:} A) 0\%; B) 30\%; C) 50\%; D) 70\%
\\ \textit{Note:} The percent chance that it will NOT rain is calculated by subtracting the probability of rain (70\%) from 100\%, resulting in a 30\% chance of no rain.
\item \textbf{Answer:} D
\\ \textit{Options:} A) 3; B) 5; C) 20; D) 60
\\ \textit{Note:} Three fifths of 100 is calculated by multiplying 100 by the fraction 3/5, which equals 60.
\item \textbf{Answer:} D
\\ \textit{Options:} A) 8; B) 84/3; C) 64\ensuremath{\sqrt{2}}/3; D) 104/3
\\ \textit{Note:} The correct answer is 104/3 because the area between the parabola y = $x^2$ and the horizontal lines y = 1 and y = 9 can be calculated by integrating the difference between the line equations and the parabola over the interval determined by their intersection points, yielding the result of 104/3.
\item \textbf{Answer:} A
\\ \textit{Options:} A) -7; B) 5; C) 7; D) 27
\\ \textit{Note:} The correct answer is -7 because when substituting n = 11 into the expression 10 - (n + 6), it simplifies to 10 - (11 + 6) = 10 - 17 = -7.
\end{enumerate}

\section*{True / False}
\begin{enumerate}[label=\arabic*.]
\setcounter{enumi}{5}
\item \textbf{Answer:} True
\\ \textit{Options:} A) True; B) False
\\ \textit{Note:} The statement is true because the function y = (ln x)/x reaches its maximum at x = e, where its value is indeed 1/e, as confirmed by finding the first derivative and analyzing critical points.
\item \textbf{Answer:} True
\\ \textit{Options:} A) True; B) False
\\ \textit{Note:} When you divide 1,248 by 7, it equals 178 with a remainder of 2, since 7 times 178 is 1,246, and subtracting this from 1,248 leaves a remainder of 2.
\item \textbf{Answer:} False
\\ \textit{Options:} A) True; B) False
\\ \textit{Note:} The point (5, 5) cannot lie on the circle that passes through (3, 4) and (5, 7) because it does not satisfy the equation of the circle formed by those two points, which requires the third point to be equidistant from the center of the circle defined by the other two points.
\item \textbf{Answer:} True
\\ \textit{Options:} A) True; B) False
\\ \textit{Note:} The statement is true because it accurately describes Noah's biking speed as 10 miles per hour, which is a valid measure of speed in the context of distance and time.
\item \textbf{Answer:} False
\\ \textit{Options:} A) True; B) False
\\ \textit{Note:} The P-value cannot exceed 1, as it represents the probability of obtaining a test statistic at least as extreme as the one observed under the null hypothesis, thus making the statement false.
\end{enumerate}

\section*{Long Form Response}
\begin{enumerate}[label=\arabic*.]
\setcounter{enumi}{10}
\item \textbf{Answer:} To pass the algebra test, you need to score at least 80\% out of 35 problems. To find out how many problems you can miss, first calculate 80\% of 35, which is 28 problems. This means you must answer at least 28 problems correctly to pass. Since there are 35 problems in total, you can miss 35 - 28 = 7 problems while still achieving a passing score. Therefore, you can miss up to 7 problems and still pass the test.
\\ \textit{Note:} correct because it accurately calculates that achieving at least 80\% of 35 problems requires answering 28 correctly, allowing for a maximum of 7 missed problems to still pass the test.
\item \textbf{Answer:} To find the area of a square whose perimeter equals that of an equilateral triangle with a side length of 8, we first calculate the perimeter of the triangle. The perimeter of an equilateral triangle is given by multiplying the side length by 3, so for a triangle with a side length of 8, the perimeter is 8 \ensuremath{\times} 3 = 24. Since the perimeter of the square is equal to 24, we set the equation for the perimeter of a square, which is 4 times the side length (s), equal to 24: 4 s = 24. Solving for s gives us s = 6. Finally, the area of the square is calculated by squaring the side length: 6 \ensuremath{\times} 6 = 36. Thus, the area of the square is 36.
\\ \textit{Note:} correct because it accurately calculates the perimeter of the equilateral triangle, equates it to the perimeter of the square to find the side length, and uses the correct formula for both shapes to ultimately determine the area of the square.
\end{enumerate}

\vfill
\noindent\textit{End of Answer Key}
\end{document}